

%%%%%%%%%%%%%%%%%%%%%%%%%%%%%%%%%%%%%%%%%%%%%%%%%%%%%%%%%%%%%%%%%%%%%%%%%%%%%%%
%                                                                             %
%  TITEL:  Abschlussarbeit 																                    %
%          --- Hauptdokument ---                                              %
%                                                                             %
%  DATEI:  AbschlussLatex.tex                                                 %
%                                                                             %
%  DATUM:  01.07.2017                                                         %
%                                                                             %
%%%%%%%%%%%%%%%%%%%%%%%%%%%%%%%%%%%%%%%%%%%%%%%%%%%%%%%%%%%%%%%%%%%%%%%%%%%%%%%


%--- Vorspann ----------------------------------------------------------------%

\documentclass[fontsize = 12pt,					%Schriftgr��e
               paper = a4,						%Papierformat
               headings = small,					%Gr��e der �berschriften
               open=right,						%Abschnitte beginnen rechts
               cleardoublepage = empty,			%leere Seiten ohne Kopfzeile
			   BCOR = 10mm,						%Binde Korrektur
               captions = tableheading,			%Tabellen mit �berschriften
               bibliography = totoc,				%Literatur- ins Inhaltsverzeichnis
               listof = totoc,					%Verzeichnisse ins Inhaltsverzeichnis (totoc)
               twoside = false,					%doppelseitiges Layout
               american]							%Dokumentklasse
              {scrbook}							

 
%--- Kopf- und Fu�zeilen -----------------------------------------------------%

\usepackage{scrpage2}
\usepackage{scrhack}
\pagestyle{scrheadings}
\automark[section]{chapter}
\addtokomafont{pageheadfoot}{\normalfont\sffamily}
\addtokomafont{pagenumber}{\normalfont\sffamily\bfseries}
\setheadsepline{0.4pt}
\clearscrheadfoot
\cehead[]{\leftmark}
\cohead[]{\rightmark}
\ohead[]{\pagemark}
\ofoot[\pagemark]{}


%--- Tabellen- und Abbildungsbeschreibungen ----------------------------------%

\addtokomafont{caption}{\small}
\addtokomafont{captionlabel}{\sffamily\bfseries}


%--- Gliederungstiefe (Bezifferung/Inhaltsverzeichnis) -----------------------%

\setcounter{secnumdepth}{2}
\setcounter{tocdepth}{1}


%--- Eingabekodierung --------------------------------------------------------%

%------ Linux ----------------------------------------------------------------%
%\usepackage[latin1]{inputenc}

%------ Windows --------------------------------------------------------------%
\usepackage[ansinew]{inputenc}


%--- Schrift -----------------------------------------------------------------%
\usepackage{lmodern}
\usepackage[T1]{fontenc}

% Spezielles Seitenlayout:
\usepackage[Sonny]{fncychap}		%Sonny/Lenny/Glenn/Conny/Rejne/Bjarne
\ChTitleVar{\LARGE\bfseries}
\ChNameVar{\large\bfseries}

%--- Math stuff --------------------------------------------------------------%
\usepackage{amsmath}

\usepackage{bm}
\usepackage{mathtools}
\newcommand{\myM}[1]{\bm{\mathit{#1}}} %Bold, italic, capital expression of matrices; vectors bold, italic, small
\DeclareMathOperator*{\argmin}{argmin}
\newcommand{\bfun}{\myM{f}}
\newcommand{\bl}{\myM{l}}
\newcommand{\bq}{\myM{q}}
\newcommand{\btau}{\myM{\tau}}
\newcommand{\bu}{\myM{u}}
\newcommand{\bv}{\myM{v}}
\newcommand{\bx}{\myM{x}}
\newcommand{\bQ}{\myM{Q}}
\newcommand{\bV}{\myM{V}}
\newcommand{\bdq}{\dot{\myM{q}}}	
\newcommand{\bddq}{\ddot{\myM{q}}}
\newcommand{\du}{\delta\myM{u}}
\newcommand{\dx}{\delta\myM{x}}



%--- Weitere Pakete ----------------------------------------------------------%
\usepackage{setspace}            % Zeilenabstand einstellbar
\usepackage{hyperref}
\usepackage{subcaption}
\usepackage{graphicx}
\usepackage{epstopdf}
\usepackage[table,xcdraw]{xcolor} %tabellen mit farben
\usepackage[export]{adjustbox} %f�r die gr��e von bildern
\usepackage{listings} %statt verbatim?
\usepackage{bibgerm}
\usepackage{float}        % u.a. genaue Plazierung von Gleitobjekten mit H
\usepackage{textcomp} 		%f�r \textregistred in �berschrift
\usepackage{tabularx}
\renewcommand{\bf}{\normalfont \bfseries} 
\usepackage{fancyvrb}  %zur Darstellung von Quelltexten
\usepackage[main=american,ngerman]{babel}
\usepackage[babel]{csquotes}
\usepackage{siunitx}  % degree symbol by \ang{degrees}
\usepackage[acronym,nonumberlist,nopostdot, automake]{glossaries} % acronym option f�r seperates Glossar und Abk�rzungsliste -> sonst zusammen
\clubpenalty = 10000             % Schusterjungen und Hurenkinder vermeiden
\widowpenalty = 10000            % siehe dazu 
\displaywidowpenalty = 10000     % http://www.cs.uu.nl/wais/html/na-dir/de-tex-faq/part5.html

%Literatur
\usepackage{cite}                % Sortierte und zusammengefasste Zitatnummern 
\usepackage[numbers]{natbib}     % Refernz mit nummern UND \citauthor{} m�glich 
\typearea[current]{current}      % Neuberechnung des Satzspiegels mit alten Werten nach ?nderung von Zeilenabstand,etc
\renewcommand{\arraystretch}{1.2} %Tabellenh�he vergr��ern
\usepackage{subcaption}			%For subtables (table a and table b automatically)

%\renewcommand{\thesubfigure}{\arabic{subfigure}} %makes subcaptions in figs/tbls 1,2,3 instead a,b,c
\makeglossaries					%Erstelle generell Glossare

%%% abbreviations:
\newacronym[longplural={degrees of freedom}]{DOF}{DOF}{degree of freedom}
\newacronym{COG}{COG}{center of gravity}

%%% nomenclature:
%\newglossaryentry{space_joint}{name={$\myM{q}=(q_1,q_2)^T$},description={minimal coordinates in joint space}}



		%include .tex data with definitions 
\usepackage{cleveref}
\captionsetup[subfigure]{subrefformat=simple,labelformat=simple}
%%%%%%%% Hints for glossary:
% call in text with \gls{label} or \glspl for plural (Gls/Glspl for capital)


%--- Hauptdokument -----------------------------------------------------------%

\begin{document}
\renewcommand*{\chapterpagestyle}{empty}

\frontmatter

%-----Titelblatt, Abstract, Kurzfassung, Inhaltsverzeichnis, Figuren, Tabellen,Terms (r�m. Ziffern)-----
%-----------------------------------------------------------------------------%
%                                                                             %
%    T I T E L S E I T E                                                      %
%                                                                             %
%-----------------------------------------------------------------------------%
%
\thispagestyle{empty}
\hypersetup{pageanchor=false}
%
\begin{flushleft}
University of Duisburg-Essen\\[1ex]
Faculty of Engineering\\[1ex]
Chair of Mechatronics \\[2ex]
\end{flushleft}
%
\vfill\vfill\vfill\vfill
%
\begin{center}
\Large
\textsf{\textbf{Highly-Dynamic Movements of a Humanoid Robot Using Whole-Body Trajectory Optimization}}\\[5ex]
\large
\textbf{Master Thesis}\\[1ex]
\textbf{Maschinenbau (M.Sc.)}\\[5ex]
Julian E�er\\[1ex]
Student ID: 3015459\\[5ex]
\begin{tabular}{ll}
\textbf{First examiner} & Prof. Dr. Dr. h.c. Frank Kirchner (DFKI)\\[1ex]
\textbf{Second examiner} & Dr.-Ing. Tobias Bruckmann (UDE)\\[1ex]
\textbf{Supervisor} & Dr. rer. nat. Shivesh Kumar (DFKI)\\[1ex]
\textbf{Supervisor} & Dr. Olivier Stasse (LAAS-CNRS)\\[1ex]
%\textbf{Supervisor} & Dr. Carlos Mastalli (University of Edinburgh)\\[1ex]
\end{tabular}\\[7ex]
\normalsize
September 22, 2020
\end{center}
%
\vfill\vfill\vfill\vfill\vfill\vfill
%
\begin{center}
	\includegraphics[width=5.2cm]{img/logo_ude}\hspace{20pt}
	\includegraphics[width=8cm]{img/logo_dfki.png}
\end{center}
\addchap*{Declaration}


\thispagestyle{empty}

This study was carried out at the Robotics Innovation Center of the German Research Center for Artificial Intelligence (DFKI) in the Advanced AI Team on Mechanics \& Control under the direction of Dr. Shivesh Kumar. 
\vspace*{5.2cm}\\
I declare that this thesis was composed by myself, that the work contained herein is my own except where explicitly stated otherwise  in the text, and that this work has not been submitted for any other degree or professional qualification except as specified. 

\bigskip
\bigskip 
\bigskip 
\bigskip
\bigskip

\begin{flushright}
\begin{tabular}{c}
Bremen, September 18, 2020\\
\bigskip \\\\
\hrulefill\hrulefill\hrulefill\hrulefill\hrulefill\hrulefill \\
Julian E�er\\
\hspace{6cm} \\
\end{tabular}
\end{flushright}



\thispagestyle{empty}
\chapter*{Abstract}

In order to further close the gap between robots and their natural counterparts, current research is driving towards exploiting the natural dynamics of robots. This requires a rethinking about the way we control robots for moving in a more dynamic, efficient and natural way. 

Dynamic bipedal locomotion is a challenging problem and remains an open area of research. A key characteristic is the decoupling between the center of mass and the multi-body dynamics. Particular difficulties arise from effective underactuation, the mechanism complexity, as well as nonlinear and hybrid dynamics. 

A common approach is to decompose this problem into smaller sub-problems that are solved sequentially. Many state-of-the-art frameworks rely on trajectory optimization based on reduced centroidal dynamics, which is transferred to a whole-body trajectory via feedback linearization using Inverse Kinematics (IK) or Inverse Dynamics (ID). Recent research indicates that solving this mapping with a local optimal control solver, namely Differential Dynamic Programming (DDP), instead of IK/ID, produces more efficient motions, with lower forces and impacts. 

This master thesis contributes to the research field of dynamic bipedal locomotion by applying, evaluating and extending DDP-based whole-body trajectory optimization, pursuing three objectives: 
First, we develop an approach to constrain the optimal control problem allowing DDP to produce inherently balanced motions. 
Second, we evaluate our approach for quasi-static and dynamic motions in a real-time physics simulation and in real-world experiments on the lightweight and biologically inspired RH5 humanoid robot.
Third, we examine the limits of the derived whole-body planning approach and the system design via solving highly-dynamic movements. 

\vfill\vfill
\textbf{Keywords:} Humanoid Robots, Dynamic Bipedal Walking, Motion Planning, Multi-Contact Optimal Control, Differential Dynamic Programming, Whole-Body Trajectory Optimization  
 












\thispagestyle{empty}
\chapter*{Kurzfassung}

Die Bewegungsplanung f{\"u}r Laufroboter ist eine anspruchsvolle Aufgabe und bleibt ein aktives Forschungsgebiet. Besondere Schwierigkeiten ergeben sich aus der effektiven Unteraktuierung, der Komplexit{\"a}t des Mechanismus sowie der nichtlinearen und hybriden Dynamik. Ein {\"u}blicher Ansatz besteht darin, das Problem in kleinere Teilprobleme zu zerlegen, die nacheinander gel{\"o}st werden. Neue Forschungsergebnisse zeigen, dass Algorithmen aus dem Bereich der optimalen Regelung, insbesondere Differenzielle Dynamische Programmierung (DDP), effizientere Bewegungen mit geringeren Kr{\"a}ften und St{\"o}{\ss}en erzeugen.

Diese Masterarbeit leistet einen Beitrag in diese Richtung, indem eine DDP-basierte Ganzk{\"o}rper-Trajektorienoptimierung angewendet, bewertet und erweitert wird, wobei drei Ziele verfolgt werden. Erstens entwickeln wir eine Methode zur Beschr{\"a}nkung von DDP-artigen Algorithmen, um stabile Bewegungspl{\"a}ne zu erzeugen. Zweitens wird der vorgeschlagene Ansatz zur Bewegungsplanung f{\"u}r quasi-statische und dynamische Bewegungen in einer physikalischen Echtzeitsimulation und in realen Experimenten mit dem leichten und biologisch inspirierten humanoiden Roboter RH5 evaluiert. Schlie{\ss}lich werden die Grenzen des Ansatzes und des Systemdesigns anhand hochdynamischer Bewegungen untersucht.

\vfill
\noindent\textbf{Stichworte:} Bewegungsplanung, Differenzielle Dynamische Programmierung , Dynamisches Bipedales Laufen, Ganzk{\"o}rper-Trajektorienoptimierung, Humanoide Roboter, Optimale Regelung

\printglossary[title=List of Symbols]
\printglossary[type=\acronymtype,title=Acronyms]		         
\tableofcontents             
\listoffigures 
\listoftables 


%-----Hauptkapitel-----
\mainmatter             % numerische Ziffern
\clearpage
\pagestyle{scrheadings}  % Kopfzeilen
\thispagestyle{scrheadings}  % Kopfzeilen
%%-----------------------------------------------------------------------------%
%                                                                             %
%    K A P I T E L   1                                                        %
%                                                                             %
%-----------------------------------------------------------------------------%

\chapter{Introduction}\label{c1}

\section{Motivation}
\section{Related Work}
%\subsection{Model Induced Design} (Passive Dynamic Walking, Template Models)
%\subsection{Traditional Legged Locomotion Planning}
%\subsection{Trajectory Optimization} (Sequential, Simultaneous OR Direct, Indirect, DDP)
%\subsection{Series-Parallel Hybrid Robots}


\section{Contribution and Structure}


\section{RH5 Humanoid Robot}
The derived approaches for constrained \gls{DDP} have been tested both in simulation and real-world experiments on a full-size humanoid robot. RH5 is a lightweight and biologically inspired humanoid that has recently been developed at DFKI Robotics Innovation Center\cite{peters2017konstruktion}.

The RH5 humanoid robot (see \cref{img:rh5_robot}) is designed to mimic the human anatomy with a total size of 200cm, a weight of 62kg and a total of 32 \gls{DoF}. The two legs account for 12 \gls{DoF}, the torso and neck kinematics each for three and the arms and grippers of the robot for 16 \gls{DoF}. In order to achieve a high dynamic performance, the robot's design follows a series-parallel hybrid approach. Consequently,linkages and parallel mechanisms are utilized in most of the robots joints, e.g. the hip-flexion-extension, knee, ankle, torso and wrist. A comparison of RH5 with other state of the art humanoid robots revealed several advantages of this design approach, including better maximum velocity and torque of the ankle as well as an advantageous weight of the lower leg \cite{kumar2020survey}. The interested reader can find a comprehensive introduction on series-parallel hybrid robots in \cite[Ch.2]{kumar2019modular}. 

\begin{figure}[h!]
\centering	
\includegraphics[width=.3\textwidth]{img/rh5_robot.png}
\caption{The recently presented RH5 is a lightweight and biologically inspired humanoid used as experimental platform within this thesis.}
\label{img:rh5_robot}
\end{figure} 

% Add to Appendix: 
% - Gelenkwinkel
% - Gelenkachsen Darstellung
% - URDF (abstract and full)

%%-----------------------------------------------------------------------------%
%                                                                             %
%    K A P I T E L   3                                                      %
%                                                                             %
%-----------------------------------------------------------------------------%

\chapter{Contact Stability Constrained DDP}\label{c3}
\section{Idea: Constraining the CoP of each Contact Surface}

\section{Center of Pressure Constraints}
\subsection{Center of Pressure (CoP) Conditions}
Conditions for \gls{CoP} inside convex hull:
\begin{align}
\begin{split}
Pitch:& -X \leq C_x \leq X \\
Roll:& -Y \leq C_y \leq Y
\end{split}
\end{align}
Represented by the four inequality conditions:
\begin{align}
\begin{split}
-X &\leq C_x \\
C_x &\leq X \\
-Y &\leq C_y \\
C_y &\leq Y \\
\end{split}
\end{align}

\subsection{CoP Computation}
\begin{equation}
\myM{p}_{CoP} = \dfrac{\myM{n}\times\myM{M}}{\myM{F}\cdot\myM{n}} = \dfrac{\begin{bmatrix} n_x \\ n_y \\ n_z \end{bmatrix} \times \begin{bmatrix} t_x \\ t_y \\ t_z \end{bmatrix}}{\begin{bmatrix} f_x \\ f_y \\ f_z \end{bmatrix} \cdot \begin{bmatrix} n_x \\ n_y \\ n_z \end{bmatrix}} = 
\begin{bmatrix} n_yt_z - n_zt_y \\ n_zt_x-n_xt_z \\ n_xt_y-n_yt_x \end{bmatrix}\cdot \dfrac{1}{f_xn_x+f_yn_y+f_zn_y}
\end{equation}

\subsection{CoP Constraints: General Case}
These conditions are be implemented via: 
\begin{equation}
\begin{bmatrix} -Yn_0 & -Yn_1 & -Yn_2 & -n_2 & 0 & n_0 \\
-Yn_0 & -Yn_1 & -Yn_2 & n_2 & 0 & -n_0 \\
-Xn_0 & -Xn_1 & -Xn_2 & 0 & n_2 & -n_1 \\
-Xn_0 & -Xn_1 & -Xn_2 & 0 & -n_2 & n_1 \\ \end{bmatrix} \cdot
\begin{bmatrix} f^x \\ f^y \\ f^z \\ \tau^x \\ \tau^y \\ \tau^z \end{bmatrix} \leq
\begin{bmatrix} 0 \\ 0 \\ 0 \\ 0 \end{bmatrix}
\end{equation}

\subsection{CoP Constraints: Special Case of Horizontal Floor}
For the case of horizontal floor, with the according normal vector $\myM{n}=[0,0,1]$, the \gls{CoP} can be computed to: 
\begin{equation}
\myM{p}_{CoP} = \begin{bmatrix} -t_y/f_z \\ t_x/f_z \\ 0 \end{bmatrix}
\end{equation}
Represented by the four inequality conditions:
\begin{align}
\begin{split}
\tau^x &\leq Yf^z \\
-\tau^x &\leq Yf^z \\
\tau^y &\leq Yf^z \\
-\tau^y &\leq Yf^z 
\end{split}
\end{align}
Briefly summarized as:
\begin{align}
\begin{split}
\mid\mid\tau^x\mid\mid &\leq Yf^z \\
\mid\mid\tau^y\mid\mid &\leq Xf^z 
\end{split}
\end{align}
These conditions are be implemented via: 
\begin{equation}
\begin{bmatrix} 0 & 0 & -Y & 1 & 0 & 0 \\
0 & 0 & -Y & -1 & 0 & 0 \\
0 & 0 & -X & 0 & 1 & 0 \\
0 & 0 & -X & 0 & -1 & 0 \end{bmatrix} \cdot
\begin{bmatrix} f^x \\ f^y \\ f^z \\ \tau^x \\ \tau^y \\ \tau^z \end{bmatrix} \leq
\begin{bmatrix} 0 \\ 0 \\ 0 \\ 0 \end{bmatrix}
\end{equation}





\subsection{Backup: Friction Cone constraints}
\begin{align}
\begin{split}
\mid\mid f^x\mid\mid &\leq \mu f^z \\
\mid\mid f^y\mid\mid &\leq \mu f^z \\
f^z &> 0
\end{split}
\end{align}
For the case of four edges of the linear approximation of the friction cone, the equations become:
\begin{equation}
\begin{bmatrix} 1 & 0 & -\mu \\
-1 & 0 & -\mu \\
0 & 1 & -\mu \\
0 & -1 & -\mu \\
0 & 0 & -\mu \\ \end{bmatrix} \cdot
\begin{bmatrix} f^x \\ f^y \\ f^z \end{bmatrix} \leq
\begin{bmatrix} 0 \\ 0 \\ 0 \\ 0 \\ 0 \end{bmatrix}
\end{equation}

\section{Integration into the Crocoddyl Framework}
\subsection{Inequality Constraints by Penalization}
%CostModelSum, Residual, Activation
\subsection{Computing the Residual}
\subsection{Unittesting Numerical Differentiation} 





%%-----------------------------------------------------------------------------%
%                                                                             %
%    K A P I T E L   4                                                        %
%                                                                             %
%-----------------------------------------------------------------------------%

\chapter{Experimental Validation}\label{c4}




%%-----------------------------------------------------------------------------%
%                                                                             %
%    K A P I T E L   5                                                        %
%                                                                             %
%-----------------------------------------------------------------------------%

\chapter{Highly-Dynamic Movements}\label{c5}

\section{Formulation of the Optimization Problem}
\section{Simulation Results for Increasing Task Complexity}
\subsection{A Simple Vertical Jump}
\subsection{Forward Jumping Over Multiple Obstacles}
\section{Identification of Limits in System Design}
\subsection{Analysis of Joint Limits}
\subsection{Analysis of Torque Limits}




\chapter{Introduction}\label{c1}
\section{Motivation}
\section{Related Work}
%\subsection{Model Induced Design} (Passive Dynamic Walking, Template Models)
%\subsection{Traditional Legged Locomotion Planning}
%\subsection{Trajectory Optimization} (Sequential, Simultaneous OR Direct, Indirect, DDP)
%\subsection{Series-Parallel Hybrid Robots}
\section{Contribution}
\section{Structure}

%-----------------------------------------------------------------------------%
%                                                                             %
%    K A P I T E L   2                                                        %
%                                                                             %
%-----------------------------------------------------------------------------%

\chapter{Background: Dynamics of Legged Robots}\label{c2}






\chapter{Dynamic Bipedal Walking}
\section{Formulation of the Optimization Problem}
\section{Inequality Constraints for Physical Compliance}
\section{Trajectories for Increasing Mechanism Complexity}

\chapter{Highly-Dynamic Movements}
\section{Formulation of the Optimization Problems}
\section{Trajectories for Increasing Task Complexity}
\section{Identification of Limits in System Design}

\chapter{Experimental Validation}
\section{Validation in Real-Time Physics Simulation}
\section{Validation on the RH5 Humanoid Robot}

\chapter{Conclusion and Outlook}
\section{Summary}
\section{Future Directions}



%-----Anhang-----
%-----------------------------------------------------------------------------%
%                                                                             %
%    Anh�nge					                                             					  %
%                                                                             %
%-----------------------------------------------------------------------------%

\appendix
\chapter{Appendix}

\subsection{Carlos Talk: Essentials}
In the end we want to solve a bilevel (nested) optimization
\begin{equation} 
\begin{split}
\myM{X}^*,\myM{U}^* = \arg\min_{\myM{X},\myM{U}} \sum_{k=0}^{N-1} task(x_k,u_k) \\
x_k=\arg\min physics(x_k,u_k), \\
s.t. \quad constraints(x_k,u_k)
\end{split}
\end{equation}

Which more formally looks like
\begin{equation*} 
\begin{split}
\myM{X}^*,\myM{U}^*=
\begin{Bmatrix} \myM{x}^*_0,\cdots,\myM{x}^*_N \\
      \myM{u}^*_0,\cdots,\myM{u}^*_N
\end{Bmatrix} =
\arg\min_{\myM{X},\myM{U}} l_N(x_N)+\sum_{k=0}^{N-1} \int_{t_k}^{t_k+\Delta t} l(\myM{x},\myM{u})dt \\
s.t.\quad\myM{\dot{v}}, \myM{\lambda}=\arg\min_{\myM{\dot{v}}, \myM{\lambda}}||\myM{\dot{v}}-\myM{\dot{v}}_{free}||_M, \\
\myM{x}\in \mathcal{X}, \myM{u}\in \mathcal{U}
\end{split}
\end{equation*}

KKT Matrix:
\begin{equation*} 
\begin{split}
\myM{X}^*,\myM{U}^* = \arg\min_{\myM{X},\myM{U}} \sum_{k=0}^{N-1} task(x_k,u_k) \\
KKT-Dynamics(x_k,u_k) 
\end{split}
\end{equation*}

Multi-contact dynamics as holonomic constraints:
\begin{equation*} 
\left[\begin{matrix}
  \myM{\dot{v}} \\
  -\myM{\lambda} \\
\end{matrix}\right] 
=
\left[\begin{matrix}
 \myM{M} & \myM{J}^{\top}_c \\
 {\myM{J}_{c}} & \myM{0} \\
\end{matrix}\right]^{-1}
\left[\begin{matrix}
 \myM{\tau}_b \\ -\myM{a_0}
\end{matrix}\right]
\end{equation*}


\section{Crocoddyl: Contact RObot COntrol by Differential DYnamic programming Library (Wiki Home)}
\subsection{Welcome to Crocoddyl}
Crocoddyl is an optimal control library for robot control under contact sequence. Its solver is based on an efficient Differential Dynamic Programming (DDP) algorithm. Crocoddyl computes optimal trajectories along to optimal feedback gains. It uses Pinocchio for fast computation of robot dynamics and its analytical derivatives.
Crocoddyl is focused on multi-contact optimal control problem (MCOP) which as the form:
$$\myM{X}^*,\myM{U}^*= \begin{Bmatrix} \myM{x}^*_0,\cdots,\myM{x}^*_N \\ \myM{u}^*_0,\cdots,\myM{u}^*_N \end{Bmatrix} = \arg\min_{\myM{X},\myM{U}} \sum_{k=1}^N \int_{t_k}^{t_k+\Delta t} l(\myM{x},\myM{u})dt$$
subject to
$$\myM{\dot{x}} = \myM{f}(\myM{x},\myM{u}),$$
$$\myM{x}\in\mathcal{X}, \myM{u}\in\mathcal{U}, \boldsymbol{\lambda}\in\mathcal{K}.$$
where
\begin{itemize}
\item the state $\myM{x}=(\myM{q},\myM{v})$ lies in a manifold, e.g. Lie manifold $\myM{q}\in SE(3)\times \myM{R}^{n_j}$,
\item the system has underactuacted dynamics, i.e. $\myM{u}=(\myM{0},\boldsymbol{\tau})$,
\item $\mathcal{X}$, $\mathcal{U}$ are the state and control admissible sets, and
\item $\mathcal{K}$ represents the contact constraints.
\end{itemize}

Note that $\boldsymbol{\lambda}=\myM{g}(\myM{x},\myM{u})$ denotes the contact force, and is dependent on the state and control.

Let's start by understanding the concept behind crocoddyl design.

\subsection{Action Models}
In crocoddyl, an action model combines dynamics and cost models. Each node, in our optimal control problem, is described through an action model. Every time that we want describe a problem, we need to provide ways of computing the dynamics, cost functions and their derivatives. All these is described inside the action model.
To understand the mathematical aspects behind an action model, let's first get a locally linearize version of our optimal control problem as:
$$\myM{X}^*(\myM{x}_0),\myM{U}^*(\myM{x}_0) = \arg\min_{\myM{X},\myM{U}} = cost_T(\delta\myM{x}_N) + \sum_{k=1}^N cost_t(\delta\myM{x}_k, \delta\myM{u}_k)$$
subject to
$$dynamics(\delta\myM{x}_{k+1},\delta\myM{x}_k,\delta\myM{u}_k)=\myM{0},$$
where
$$cost_T(\delta\myM{x}_k) = \frac{1}{2} \begin{bmatrix} 1 \\ \delta\myM{x}_k \end{bmatrix}^\top \begin{bmatrix} 0 & \myM{l_x}^\top_k \\ \myM{l_x}_k & \myM{l_{xx}}_k \end{bmatrix} \begin{bmatrix} 1 \\ \delta\myM{x}_k \end{bmatrix},$$     

$$cost_t(\delta\myM{x}_k,\delta\myM{u}_k) = \frac{1}{2} \begin{bmatrix} 1 \\ \delta\myM{x}_k \\ \delta\myM{u}_k \end{bmatrix}^\top \begin{bmatrix} 0 & \myM{l_x}^\top_k & \myM{l_u}^\top_k\\ \myM{l_x}_k & \myM{l_{xx}}_k & \myM{l_{ux}}^\top_k\\ \myM{l_u}_k & \myM{l_{ux}}_k & \myM{l_{uu}}_k \end{bmatrix} \begin{bmatrix} 1 \\ \delta\myM{x}_k \\ \delta\myM{u}_k \end{bmatrix}$$

$$dynamics(\delta\myM{x}_{k+1},\delta\myM{x}_k,\delta\myM{u}_k) = \delta\myM{x}_{k+1} - (\myM{f_x}_k\delta\myM{x}_k + \myM{f_u}_k\delta\myM{u}_k)$$

\subsubsection{Notes}
\begin{itemize}
\item An action model describes the dynamics and cost functions for a node in our optimal control problem.
\item Action models lie in the discrete time space.
\item For debugging and prototyping, we have also implemented NumDiff abstractions. These computations depend only in the defining of the dynamics equation and cost functions. However to asses efficiency, crocoddyl uses analytical derivatives computed from Pinocchio.
\end{itemize}

\subsubsection{Differential and Integrated Action Models}

It's often convenient to implement action models in continuous time. In crocoddyl, this continuous-time action models are called Differential Action Model (DAM). And together with predefined Integrated Action Models (IAM), it possible to retrieve the time-discrete action model needed by the solver.
At the moment, we have the following integration rules:
\begin{itemize}
\item simpletic Euler and
\item Runge-Kutta 4.
\end{itemize}

\subsubsection{Add On from Introduction.jpnb}
Optimal control solvers often need to compute a quadratic approximation of the action model (as previously described); this provides a search direction (computeDirection). Then it's needed to try the step along this direction (tryStep).

Typically calc and calcDiff do the precomputations that are required before computeDirection and tryStep respectively (inside the solver). These functions update the information of:
\begin{itemize}
\item \textbf{calc}: update the next state and its cost value
 $$\delta\myM{\dot{x}}_{k+1} = \myM{f}(\delta\myM{x}_k,\myM{u}_k)$$
\item \textbf{calcDiff}: update the derivatives of the dynamics and cost (quadratic approximation)
 $$\myM{f_x}, \myM{f_u} \hspace{1em} (dynamics)$$
 $$\myM{l_x}, \myM{l_u}, \myM{l_{xx}}, \myM{l_{ux}}, \myM{l_{uu}} \hspace{1em} (cost)$$
\end{itemize}


\subsection{State and its Integrate and Difference Rules}
General speaking, the system's state can lie in a manifold $M$ where the state rate of change lies in its tangent space $T_\mathbf{x} M$. There are few operators that needs to be defined for different routines inside our solvers:
$$\myM{x}_{k+1} = integrate(\myM{x}_k,\delta\myM{x}_k) = \myM{x}_k \oplus \delta\myM{x}_k$$
$$\delta\myM{x}_k = difference(\myM{x}_{k+1},\myM{x}_k) = \myM{x}_{k+1} \ominus \myM{x}_k$$
where $\mathbf{x}\in M$ and $\delta\mathbf{x}\in T_\mathbf{x} M$.
And we also need to defined the Jacobians of these operators with respect to the first and second arguments:
$$\frac{\partial \myM{x}\oplus\delta\myM{x}}{\partial \myM{x}}, \frac{\partial \myM{x}\oplus\delta\myM{x}}{\partial\delta\myM{x}} =Jintegrante(\myM{x},\delta\myM{x})$$
$$\frac{\partial\myM{x}_2\ominus\myM{x}_2}{\partial \myM{x}_1}, \frac{\partial \myM{x}_2\ominus\myM{x}_1}{\partial\myM{x}_1} =Jdifference(\myM{x}_2,\myM{x}_1)$$
For instance, a state that lies in the Euclidean space will the typical operators:
$$integrate(\myM{x},\delta\myM{x}) = \myM{x} + \delta\myM{x}$$
$$difference(\myM{x}_2,\myM{x}_1) = \myM{x}_2 - \myM{x}_1$$
$$Jintegrate(\cdot,\cdot) = Jdifference(\cdot,\cdot) = \myM{I}$$
All these functions are encapsulate inside the State class. For Pinocchio models, we have implemented the StateMultibody class which can be used for any robot model.

\section{Crocoddyl Wiki: Differential Action Model for Floating in Contact Systems (DAMFIC)}
\subsection{System Dynamics}
As you might know, a differential action model describes the systems dynamics and cost function in continuous-time. For multi-contact locomotion, we account for the rigid contact by applying the Gauss principle over holonomic constraints in a set of predefined contact placements, i.e.:
$$\begin{aligned} & \dot{\myM{v}} = \underset{\myM{a}}{\arg\min} & & \frac{1}{2}\,\|\dot{\myM{v}}-\dot{\myM{v}}_{free}\|_{\myM{M}} \\ & \textrm{subject to} & & \myM{J}_{c} \dot{\myM{v}} + \dot{\myM{J}}_c \myM{v} = \myM{0}, \end{aligned}$$

This is equality-constrained quadratic problem with an analytical solution of the form:
$$\left[\begin{matrix}\myM{M} & \myM{J}^{\top}_c \\{\myM{J}_{c}} & \myM{0}\end{matrix}\right] \left[\begin{matrix} \dot{\myM{v}} \\ -\boldsymbol{\lambda} \end{matrix}\right] = \left[\begin{matrix} \boldsymbol{\tau}_b \\ -\dot{\myM{J}}_c \myM{v}\end{matrix}\right]$$
in which
$$(\dot{\myM{v}},\boldsymbol{\lambda})\in(\myM{R}^{nv},\myM{R}^{nf})$$
are the primal and dual solutions,
$$\myM{M}\in\myM{R}^{nv\times nv}$$ is formally the metric tensor over the configuration manifold $\myM{q}\in\myM{R}^{nq}$,
$$\myM{J}_{c}= \begin{bmatrix} \myM{J}_{c_1} & \cdots & \myM{J}_{c_f}\end{bmatrix}\in\myM{R}^{nf\times nv}$$ is a stack of $f$ contact Jacobians,
$\boldsymbol{\tau}_b=\myM{S}\boldsymbol{\tau}-\myM{b}\in\myM{R}^{nv}$ is the force-bias vector that accounts for the control $\boldsymbol{\tau}\in\myM{R}^{nu}$, the Coriolis and gravitational effects $\myM{b}$, and $\myM{S}$ is the selection matrix of the actuated joint coordinates, and
$nq$, $nv$, $nu$ and $nf$ are the number of coordinates used to describe the configuration manifold, its tangent-space dimension, control commands and contact forces, respectively.

And this equality-constrained forward dynamics can be formulated using state space representation, i.e.:
$$\dot{\myM{x}}=\myM{f(x,u)}$$
where $\myM{x}=(\myM{q,v})\in\myM{R}^{nq+nv}$ and $\myM{u}=\boldsymbol{\tau}\in\myM{R}^{nu}$ are the state and control vectors, respectively. Note that $\dot{\myM{x}}$ lies in the tangent-space of $\myM{x}$, and their dimension are not the same.

\subsection{Add On from Introduction.jpnb}
\subsection{Solving the Optimal Control Problem}
Our optimal control solver interacts with a defined ShootingProblem. A \textbf{shooting problem} represents a \textbf{stack of action models} in which an action model defines a specific node along the OC problem.

First we need to create an action model from DifferentialFwdDynamics. We use it for building terminal and running action models. In this example, we employ an simpletic Euler integration rule.

Next we define the set of cost functions for this problem. One could formulate
\begin{itemize}
\item Running costs (related to individual states)
\item Terminal costs (related to the final state)
\end{itemize}
in order to penalize, for example, the state error, control error, or end-effector pose error. 

Onces we have defined our shooting problem, we create a DDP solver object and pass some callback functions for analysing its performance.
\subsubsection{Application to Bipedal Walking}
In crocoddyl, we can describe the multi-contact dynamics through holonomic constraints for the support legs. From the Gauss principle, we have derived the model as:
$$
\left[\begin{matrix}
 \myM{M} & \myM{J}^{\top}_c \\
 {\myM{J}_{c}} & \myM{0} \\
\end{matrix}\right]
\left[\begin{matrix}
 \dot{\myM{v}} \\ -\boldsymbol{\lambda}
\end{matrix}\right]
 = 
\left[\begin{matrix}
  \boldsymbol{\tau} - \myM{h} \\
  -\dot{\myM{J}}_c \myM{v} \\
\end{matrix}\right]$$.

This DAM is defined in "DifferentialActionModelFloatingInContact" class.
Given a predefined contact sequence and timings, we build per each phase a specific multi-contact dynamics. Indeed we need to describe \textbf{multi-phase optimal control problem}. One can formulate the multi-contact optimal control problem (MCOP) as follows:

$$\myM{X}^*,\myM{U}^*=
\begin{Bmatrix} \myM{x}^*_0,\cdots,\myM{x}^*_N \\
				  \myM{u}^*_0,\cdots,\myM{u}^*_N
\end{Bmatrix} =
\arg\min_{\myM{X},\myM{U}} \sum_{p=0}^P \sum_{k=1}^{N(p)} \int_{t_k}^{t_k+\Delta t} l_p(\myM{x},\myM{u})dt$$
subject to
$$ \myM{\dot{x}} = \myM{f}_p(\myM{x},\myM{u}), \text{for } t \in [\tau_p,\tau_{p+1}]$$

$$ \myM{g}(\myM{v}^{p+1},\myM{v}^p) = \myM{0}$$

$$ \myM{x}\in\mathcal{X}_p, \myM{u}\in\mathcal{U}_p, \boldsymbol{\lambda}\in\mathcal{K}_p.$$

where $\myM{g}(\cdot,\cdot,\cdot)$ describes the contact dynamics, and they represents terminal constraints in each walking phase. In this example we use the following \textbf{impact model}:

$$\myM{M}(\myM{v}_{next}-\myM{v}) = \myM{J}_{impulse}^T$$

$$\myM{J}_{impulse} \myM{v}_{next} = \myM{0}$$

$$\myM{J}_{c} \myM{v}_{next} = \myM{J}_{c} \myM{v}$$


%%% style is changed only once, need to reopen texstudio for change %%%
%\bibliographystyle{plain} %Verzeichnis nach Autor sortiert, Referenzen numerisch
%\bibliographystyle{abbrv} %Verzeichnis nach Autor sortiert, Vornamen abgek�rzt 
%\bibliographystyle{alpha} %Verzeichnis nach Autor sortiert, Referenzen aus Autorenk�rzel
\bibliographystyle{unsrtnat} %Verzeichnis nach Zitierreihenfolge sortiert und kompatibel mit natbib package
\bibliography{tex/_bibfile}	%refer to .bib file in folder
%\nocite{*}				 % Alle Quelleneintr�ge anzeigen, auch wenn sie nicht im Text referenziert sind
\end{document}
